

\medskip
\noindent\SGAInlineRefTarget{sga3:iii:proposition:4-36}\textbf{Proposition 4.36.}
\textit{Under the hypotheses of \SGARefLink{sga3:iii:theorem:4-21}{4.21}, suppose in addition that
\(Y_{\mathscr J}\) is smooth over \(S_{\mathscr J}\) and that \(S_0\)
is affine. The set of sub-\(S\)-groups \(Y\) of \(X\), flat (or smooth)
over \(S\), reducing to \(Y_{\mathscr J}\), modulo conjugation by
sections of \(X\) over \(S\) inducing the unit section of
\(X_{\mathscr J}\), is either empty or a principal homogeneous set under
the group}
\[
 H^1\!\left(
 Y_0,
 \bigl[\operatorname{Lie}(X_0/S_0)/\operatorname{Lie}(Y_0/S_0)\bigr]
 \otimes_{\mathscr O_{S_0}}\mathscr J
 \right).
\]

It is enough to verify that \SGARefLink{sga3:iii:corollary:4-29}{Corollary 4.29} applies, that is, that
\[
 d^0:
 \operatorname{Hom}_{\mathscr O_{S_0}}
   (\omega^1_{X_0/S_0},\mathscr J)
 \longrightarrow
 \operatorname{Hom}_{\mathscr O_{S_0}}
   (\mathbf n_{Y_0/X_0},\mathscr J)
\]
is surjective. This follows from the fact that the sequence \((+)\) of
\SGARefLink{sga3:iii:lemma:4-25:item:ii}{4.25(ii)} is split, since \(S_0\) is affine.

Let us finally state a corollary common to \SGARefLink{sga3:iii:theorem:4-21}{4.21} and \SGARefLink{sga3:iii:proposition:4-36}{4.36}. In fact, this
will be the only form in which we shall subsequently use the general
results of this section.

\medskip
\noindent\SGAInlineRefTarget{sga3:iii:corollary:4-37}\textbf{Corollary 4.37.}
\textit{Let \(S\) be a scheme and let \(S_0\) be the closed subscheme
defined by a nilpotent ideal \(\mathscr I\). Let \(X\) be an \(S\)-group
smooth over \(S\), and let \(Y_0\) be a sub-\(S_0\)-group of \(X_0\),
flat over \(S_0\).}
\begin{enumerate}
\item[\SGAInlineRefTarget{sga3:iii:corollary:4-37:item:i}(i)] \textit{If \(S_0\) is affine, \(Y_0\) is smooth over \(S_0\),
and if}
\[
 H^1\!\left(
 Y_0,
 \bigl[\operatorname{Lie}(X_0/S_0)/\operatorname{Lie}(Y_0/S_0)\bigr]
 \otimes_{\mathscr O_{S_0}}
 \mathscr I^{n+1}/\mathscr I^{n+2}
 \right)=0
\]
\textit{for every \(n\geq 0\), then any two sub-\(S\)-groups of \(X\),
flat (or smooth) over \(S\) and reducing to \(Y_0\), are conjugate by a
section of \(X\) over \(S\) inducing the unit section of \(X_0\).}

\item[\SGAInlineRefTarget{sga3:iii:corollary:4-37:item:ii}(ii)] \textit{If \(Y_0\) is affine and of finite presentation,
and if}
\[
 H^2\!\left(
 Y_0,
 \mathbf n_{Y_0/X_0}^{\vee}
 \otimes_{\mathscr O_{S_0}}
 \mathscr I^{n+1}/\mathscr I^{n+2}
 \right)=0
\]
\textit{for every \(n\geq 0\), then there exists a sub-\(S\)-group of
\(X\), flat over \(S\), reducing to \(Y_0\).}
\end{enumerate}

\medskip
\noindent\SGAInlineRefTarget{sga3:iii:paragraph:4-38}\textbf{4.38.}
It remains to relate the three constructions made in this Expos\'e. To
avoid inessential complications, we shall place ourselves in the
following situation: \(S_0\) is the spectrum of a field \(k\), \(S\) is
the spectrum of the dual numbers \(D(k)\), \(G\) is an \(S\)-group
smooth over \(S\), and \(K\) is a sub-\(S\)-group, smooth over \(S\) and
affine.

Put
\[
 \mathfrak g_0=\operatorname{Lie}G_0
 =\Gamma(S_0,\operatorname{Lie}G_0)
 =\operatorname{Lie}(G_0/S_0)(S_0),
 \qquad
 \mathfrak k_0=\operatorname{Lie}K_0.
\]
We have an exact sequence of \(k\)-vector spaces
\[
 0\longrightarrow\mathfrak k_0
 \xrightarrow{\ i\ }\mathfrak g_0
 \xrightarrow{\ d\ }\mathbf n^{\vee}_{K_0/G_0}
 \longrightarrow 0,
\]
which gives rise to an exact cohomology sequence
\[
\begin{aligned}
0&\longrightarrow H^0(K_0,\mathfrak k_0)
 \xrightarrow{\ i^0\ }H^0(K_0,\mathfrak g_0)
 \xrightarrow{\ d^0\ }H^0(K_0,\mathfrak g_0/\mathfrak k_0)\\
&\xrightarrow{\ \partial^0\ }H^1(K_0,\mathfrak k_0)
 \xrightarrow{\ i^1\ }H^1(K_0,\mathfrak g_0)
 \xrightarrow{\ d^1\ }H^1(K_0,\mathbf n^{\vee}_{K_0/G_0})
 \xrightarrow{\ \partial^1\ }H^2(K_0,\mathfrak k_0).
\end{aligned}
\]

These various groups all have a geometric meaning.

\begin{enumerate}
\item[\textbf{a)}]
\(H^0(K_0,\mathfrak k_0)
=\operatorname{Lie}\underline{\operatorname{Centr}}(K_0)\)
by II, \SGARefLink{sga3:ii:theorem:5-2-3}{5.2.3}.

\item[\textbf{b)}]
\(H^0(K_0,\mathfrak g_0)
=\operatorname{Lie}\underline{\operatorname{Centr}}_{G_0}(K_0)\)
 (idem).

\item[\textbf{c)}]
\(H^0(K_0,\mathfrak g_0/\mathfrak k_0)
=\operatorname{Lie}\underline{\operatorname{Norm}}_{G_0}(K_0)
/\mathfrak k_0\) (idem).

\item[\textbf{d)}]
\[
 H^1(K_0,\mathfrak k_0)
 =\operatorname{Lie}\!\left(
   \underline{\operatorname{Aut}}_{S_0\text{-gr.}}(K_0)
  \right)/\operatorname{Im}(\mathfrak k_0),
\]
where \(\operatorname{Im}(\mathfrak k_0)\) denotes the image of
\(\mathfrak k_0\) under the morphism \(\operatorname{Lie}(\operatorname{Int}_0)\)
deduced from
\[
 \operatorname{Int}_0:K_0\longrightarrow
 \underline{\operatorname{Aut}}_{S_0\text{-gr.}}(K_0).
\]
Indeed, it follows from \SGARefLink{sga3:iii:theorem:2-1:item:ii}{2.1(ii)}, applied with \(Y=X=K\) and
\(f_0=\operatorname{id}_{K_0}\), that \(Z^1(K_0,\mathfrak k_0)\) is the
group of infinitesimal automorphisms of the \(S_0\)-group \(K_0\), and
that \(H^1(K_0,\mathfrak k_0)\) is obtained by quotienting by inner
infinitesimal automorphisms, that is, by the image of \(\mathfrak k_0\).
Moreover, by II, \SGARefLink{sga3:ii:corollary:4-2-2}{4.2.2}, we also have
\[
 Z^1(K_0,\mathfrak k_0)
 =\operatorname{Lie}\!\left(
   \underline{\operatorname{Aut}}_{S_0\text{-gr.}}(K_0)
  \right).
\]

\item[\textbf{e)}]
By \SGARefLink{sga3:iii:theorem:2-1:item:ii}{2.1(ii)}, \(H^1(K_0,\mathfrak g_0)\) is the group of deviations
between homomorphisms \(K\to G\) extending the canonical immersion
\(i_0:K_0\to G_0\), modulo the deviations obtained through the action
of the inner automorphisms of \(G\) defined by elements of \(G(S)\)
which give the unit of \(G(S_0)\) (that is, elements of
\(\mathfrak g_0\)).

\item[\textbf{f)}]
By 4.36, \(H^1(K_0,\mathbf n^{\vee}_{K_0/G_0})\) is the group of
deviations between subgroups \(K'\) of \(G\) extending \(K_0\) and
flat over \(S\) (hence smooth over \(S\), cf. SGA 1, II, 4.10), modulo
the deviations obtained through the action of the inner automorphisms
of \(G\) constructed as above.

\item[\textbf{g)}]
By 3.5(ii), \(H^2(K_0,\mathfrak k_0)\) is the group of deviations
between group structures on \(K\) extending that of \(K_0\), modulo
the \(S\)-automorphisms of \(K\) inducing the identity on \(K_0\).
\end{enumerate}

We now propose to make explicit the six morphisms of the preceding
exact sequence in terms of the geometric interpretation just given.

\begin{enumerate}
\item[\textbf{1)}]
The morphisms \(i^0\) and \(d^0\) are simply those obtained by passing
to the Lie algebra (and then, for \(d^0\), to the quotient) in the
canonical monomorphisms
\[
 \underline{\operatorname{Centr}}(K_0)
 \longrightarrow
 \underline{\operatorname{Centr}}_{G_0}(K_0)
 \longrightarrow
 \underline{\operatorname{Norm}}_{G_0}(K_0).
\]
This follows immediately from the definition of the identifications
\textbf{a)}, \textbf{b)}, and \textbf{c)}.

\item[\textbf{2)}]
The morphism \(\partial^0\) is constructed as follows. Let
\[
 \bar x\in
 \operatorname{Lie}\underline{\operatorname{Norm}}_{G_0}(K_0)
 /\mathfrak k_0.
\]
Lift it to an element
\[
 x\in\operatorname{Lie}\underline{\operatorname{Norm}}_{G_0}(K_0)
 \subset\underline{\operatorname{Norm}}_G(K)(S).
\]
Then \(\operatorname{Int}(x)\) defines an automorphism of \(K\)
inducing the identity on \(K_0\), hence an element of
\(\operatorname{Lie}\underline{\operatorname{Aut}}_{S_0\text{-gr.}}(K_0)\).
Denote by \(\overline{\operatorname{Int}(x)}\) the image of this element
in
\[
 \operatorname{Lie}\underline{\operatorname{Aut}}_{S_0\text{-gr.}}(K_0)
 /\operatorname{Im}(\mathfrak k_0).
\]
Then
\begin{equation}
 \partial^0(\bar x)
 =-\overline{\operatorname{Int}(x)}
 =\overline{\operatorname{Int}(x^{-1})}.
 \tag{*}\label{sga3:iii:paragraph:4-38:equation:star}
\end{equation}

Indeed, let us compute the element of
\(\operatorname{Lie}\underline{\operatorname{Aut}}_{S_0\text{-gr.}}(K_0)\)
defined by \(\operatorname{Int}(x)\). By definition it corresponds to
an element \(a\in Z^1(K_0,\mathfrak k_0)\) such that
\[
 xyx^{-1}=a(y_0)y,
 \qquad y\in K(S'),\quad S'\longrightarrow S.
\]
But this can also be written
\[
 a(y_0)=xyx^{-1}y^{-1}=x-\operatorname{Ad}(y)x
 =-\partial(x)(y_0),
\]
whence \(a=-\partial(x)\).

On the other hand,\footnote{Editor's note: We have added the following
sentence.} the image under \(\partial\) of
\[
 x\in\operatorname{Lie}\underline{\operatorname{Norm}}_{G_0}(K_0)
 \subset\mathfrak g_0
\]
is an element of \(Z^1(K_0,\mathfrak k_0)\), whose image
\(\overline{\partial(x)}\) in \(H^1(K_0,\mathfrak k_0)\) depends only
on \(\bar x\). By definition of the connecting map \(\partial^0\),
\[
 \partial^0(\bar x)=\overline{\partial(x)}.
\]
Combined with \(a=-\partial(x)\), this proves \((*)\).

\item[\textbf{3)}]
\footnote{Editor's note: We have expanded the original in what follows,
and in \((**)\) we have corrected \(u\circ i\) to \(i\circ u\).}
Let \(i:K\to G\) denote the canonical immersion. Let
\(\bar u\in H^1(K_0,\mathfrak k_0)\) be the image of an element
\[
 u\in
 \operatorname{Lie}\underline{\operatorname{Aut}}_{S_0\text{-gr.}}(K_0)
 \subset\underline{\operatorname{Aut}}_{S\text{-gr.}}(K).
\]
Then
\begin{equation}
 i^1(\bar u)=\bar d(i,i\circ u),
 \tag{**}\label{sga3:iii:paragraph:4-38:equation:double-star}
\end{equation}
where \(\bar d(i,i\circ u)\) is the class defined in \SGARefLink{sga3:iii:remark:2-1-1}{2.1.1}.

Indeed, \(\bar u\) is the image of an element
\(v\in Z^1(K_0,\mathfrak k_0)\) such that
\(u(y)=v(y_0)y\), and \(i^1(\bar u)\) is the image in
\(H^1(K_0,\mathfrak g_0)\) of the cocycle
\(i\circ v\in Z^1(K_0,\mathfrak g_0)\). Since \(i\) is a group
morphism, the equality \(u(y)=v(y_0)y\) implies
\[
 i(u(y))=i(v(y_0))i(y).
\]
It follows that \(i\circ v=d(i,i\circ u)\), proving \((**)\).

\item[\textbf{4)}]\SGAInlineRefTarget{sga3:iii:paragraph:4-38:item:4}
Let \(i':K\to G\) be a group morphism lifting \(i_0\), let
\(\bar d(i,i')\) be the class defined in \SGARefLink{sga3:iii:remark:2-1-1}{2.1.1}, and let
\[
 d\bigl(K,i'(K)\bigr)
 \in C^1(K_0,\mathbf n^{\vee}_{K_0/G_0})
\]
be the deviation defined in \SGARefLink{sga3:iii:remark:4-5-1}{4.5.1}. By \SGARefLink{sga3:iii:theorem:4-21}{4.21},
\(d(K,i'(K))\) belongs to
\(Z^1(K_0,\mathbf n^{\vee}_{K_0/G_0})\). Denote its image in
\(H^1(K_0,\mathbf n^{\vee}_{K_0/G_0})\) by
\(\bar d(K,i'(K))\). Then, by \SGARefLink{sga3:iii:lemma:4-27:item:i}{4.27(i)},
\begin{equation}
 d^1\bigl(\bar d(i,i')\bigr)=\bar d\bigl(K,i'(K)\bigr).
 \tag{\(\dagger\)}\label{sga3:iii:paragraph:4-38:equation:dagger}
\end{equation}
\par\smallskip
\noindent\emph{Translator's note.} In these three adjacent occurrences the
French re-edition prints \(\mathbf n_{K_0/X_0}\), without the dual. The
English uses \(\mathbf n^{\vee}_{K_0/G_0}\), as required by the exact
sequence and cohomology sequence established at the start of \SGARefLink{sga3:iii:paragraph:4-38}{4.38}.

\item[\textbf{5)}]\SGAInlineRefTarget{sga3:iii:paragraph:4-38:item:5}
Finally, let \(K'\) be a subgroup of \(G\) lifting \(K_0\) and flat,
hence smooth, over \(S\). We have assumed that \(K_0\) is affine. We
therefore know that \(K\) and \(K'\) are isomorphic as schemes extending
\(K_0\) (cf. \SGARefLink{sga3:iii:proposition:0-15}{0.15}), so there exists an isomorphism of \(S\)-schemes
\[
 f:K\xrightarrow{\ \sim\ }K'
\]
inducing the identity on \(K_0\). Transport the group structure of
\(K'\) by \(f\), and let \(K_1\) be the group thereby obtained (whose
underlying scheme is thus \(K\)); that is, the group law \(\mu_1\) of
\(K_1\) is defined, for every \(S'\to S\) and every
\(x,y\in K(S')\), by
\[
 \mu_1(x,y)=f^{-1}\bigl(f(x)f(y)\bigr).
\]
By \SGARefLink{sga3:iii:remark:3-5-1}{3.5.1},\footnote{Editor's note: We have modified the original in what
follows, taking account of the additions made in \SGARefLink{sga3:iii:remark:3-5-1}{3.5.1} and \SGARefLink{sga3:iii:proposition:4-35-1}{4.35.1}.}
\(K_1\) defines a cocycle
\(\delta(K,K_1)\in Z^2(K_0,\mathfrak k_0)\) such that, for every
\(S'\to S\) and every \(x,y\in K(S')\),
\[
 \delta(K,K_1)(x_0,y_0)xy
 =\mu_1(x,y)=f^{-1}\bigl(f(x)f(y)\bigr).
\]
Then, by \SGARefLink{sga3:iii:proposition:4-35-1}{4.35.1},
\begin{equation}
 \partial^1\bar d(K,K')=\bar\delta(K,K_1).
 \tag{\(\ddagger\)}\label{sga3:iii:paragraph:4-38:equation:ddagger}
\end{equation}
\end{enumerate}

\section*{Bibliography}
\addcontentsline{toc}{section}{Bibliography}
\footnote{Editor's note: Additional references cited in this Expos\'e.}

\begin{description}
\item[{[BAlg]}] N. Bourbaki, \emph{Alg\`ebre}, Chap. I--III,
  Hermann, 1970.
\item[{[BAC]}] N. Bourbaki, \emph{Alg\`ebre commutative}, Chap. I--IV,
  Masson, 1985.
\item[{[DG70]}] M. Demazure and P. Gabriel,
  \emph{Groupes alg\'ebriques}, Masson \& North-Holland, 1970.
\item[{[Fr64]}] P. Freyd, \emph{Abelian categories}, Harper and Row,
  1964.
\end{description}

